\section{Datasets}\label{sec:datasets}
All 321 pairs are exposed or non-pristine: QuixBugs and BugsInPy are widely used benchmarks and were used in earlier evaluations of this programme, the synthetic pairs were scored by earlier-generation systems, and V7-B was shaped on them. There is no independent holdout and benchmark exposure cannot be excluded. They are development data, not evidence of field performance. The arms are (Table~\ref{tab:datasets}):
\begin{itemize}
\item \textbf{QuixBugs buggy/fixed} (31 pairs): inputs from the benchmark's declared test cases, outputs unused. Truth: divergent for the 29 pairs whose declared tests fail on the buggy and pass on the fixed version (a test-replay oracle independent of V7-B and the baselines); two are unlabelled. QuixBugs counts as one project.
\item \textbf{QuixBugs semantics-preserving negatives} (85): the fixed version against a generator-made semantics-preserving variant; label ``no observed difference''; one suspect.
\item \textbf{Synthetic changed} (78) and \textbf{synthetic equivalent} (72): generator labels; ``equivalent'' means equivalent under a looser equality than the V7-B boundary, so a witnessed pair may be a label artefact. Divergence in ``changed'' pairs may lie outside the sampled inputs.
\item \textbf{BugsInPy units} (55): real fixes, extracted single-file functions, run on synthetic battery inputs rather than project tests; nine class units are not driven. No defensible truth label exists, so all are excluded from primary denominators.
\end{itemize}
Six suspect labels are tagged and excluded but never relabelled. The primary denominators are 106 divergent-labelled and 152 no-difference-labelled pairs (258); 57 unlabelled pairs are excluded. The divergent pairs span 78 project clusters and the no-difference pairs 69; the largest cluster has 29 pairs (QuixBugs). Synthetic ``projects'' are generator base programs. No real independent project-level run exists.
\begin{table}[t]\centering\scriptsize\setlength{\tabcolsep}{3pt}
\caption{The 321 development pairs by arm and truth class (V7-B rows; all data non-pristine). Primary denominators: divergent (div.) and no-observed-difference (no-diff.) labels (258 pairs); suspect and unlabelled pairs are excluded. Source: \path{results/v7b_phase2_summary/p2_e3_20261118_092700/processed/dev_pairs_per_system.csv}.}\label{tab:datasets}
\resizebox{\linewidth}{!}{%
\begin{tabular}{lrrrrrl}
\hline
Arm & Pairs & Div. & No-diff. & Suspect & Unlab. & Truth rule \\
\hline
bugsinpy\_units & 55 & 0 & 0 & 0 & 55 & real fixes; no defensible label \\
quixbugs\_buggy\_fixed & 31 & 29 & 0 & 0 & 2 & declared tests fail on buggy, pass on fixed \\
quixbugs\_sp\_negatives & 85 & 0 & 84 & 1 & 0 & generator: semantics-preserving variant \\
synthetic\_equivalent & 72 & 0 & 68 & 4 & 0 & generator label: equivalent (looser equality) \\
synthetic\_changed & 78 & 77 & 0 & 1 & 0 & generator label: changed \\
total & 321 & 106 & 152 & 6 & 57 &  \\
\hline
\end{tabular}}
\end{table}

